\section{W[1]-hardness}\label{sec:hardness}
Theorems~\ref{thm:l0} and~\ref{thm:A} solve CEaMP in time $5^{w+1}n^{2\ell+O(1)}$, and Theorem~\ref{thm:nl} halves the exponent of $n$; all of these are XP in $\ell$. This section shows that the exponent of $n$ cannot be removed: the main question of this paper has a negative answer. The reduction is described completely and its key proofs are given here; the proofs of the long technical lemmas (the chain lemma, the profiles of the hosts, the parameter counts, the pathwidth bound) are in Appendix~\ref{app:hardness}. Four finite steps of the proof use a computer (Section~\ref{sec:hardcomp}).
\begin{theorem}[{W[1]}-hardness]\label{thm:hard}
There is a polynomial-time reduction that maps an instance $(k,n,(S_{ij}))$ of Grid Tiling with $\le$ to an instance $(G,H,\ell)$ of CEaMP with the following properties. Let $p$ be the least even integer with $p\ge k^2+5$.
\begin{enumerate}
\item[(i)] $\ell=k^2$; $H$ is a pair-maximal modification-disjoint packing; $|V(G)|=O(k^2n^2p)$.
\item[(ii)] $c(H)\le 4p+2$, and $\tw(I(H))\le\mathrm{pw}(I(H))\le 24\,k\,p^2$, independently of $n$; $I(H)$ has at most $2k$ components.
\item[(iii)] $(G,H,\ell)$ is a YES instance iff the Grid Tiling instance has a solution; if it has one, then $\OPT(G)=|H|+k^2$.
\end{enumerate}
Consequently CEaMP is W[1]-hard parameterized by $\ell+\tw(I(H))$, by $\ell+\mathrm{pw}(I(H))$ and by $\ell+c$, even for pair-maximal packings. Unless $\mathrm{FPT}=\mathrm{W[1]}$ there is no algorithm with running time $f(\ell,w)\,n^{O(1)}$, and unless ETH fails there is no algorithm with running time $f(\ell,w)\,N^{o(\sqrt{\ell})}$, where $N=|V(G)|$.
\end{theorem}



\subsection{Consequences and remaining gaps}
\begin{corollary}[The ETH bound that follows]\label{cor:hardeth}
Assume ETH. There are no computable $f$ and no algorithm that decides CEaMP in time $f(\ell,\tw(I(H)))\,N^{o(\sqrt\ell)}$ on instances with $N$ vertices.
\end{corollary}
\begin{proof}
Suppose such an algorithm runs in time $f(\ell,w)N^{h(\ell)}$ with $h(\ell)=o(\sqrt\ell)$. Apply the reduction of Theorem~\ref{thm:hard} to a Grid Tiling instance with parameter $k$: then $\ell=k^2$, $w\le 24k(k^2+6)^2=:g(k)$ and $N\le \mathrm{poly}(k)\,n^2$. The running time is at most $f(k^2,g(k))\,(\mathrm{poly}(k)n^2)^{h(k^2)}$, and $h(k^2)=o(k)$, so it is at most $f'(k)\,n^{o(k)}$ for a computable $f'$ (the factor $\mathrm{poly}(k)^{h(k^2)}=2^{o(k\log k)}$ is absorbed in $f'$). This solves Grid Tiling with $\le$ in time $f'(k)n^{o(k)}$, contradicting ETH~\cite{Cygan15}.
\end{proof}
Nothing stronger follows from this reduction: since $\ell=k^2$, the lower bound is $n^{o(\sqrt\ell)}$, while the upper bounds of Theorems~\ref{thm:A} and~\ref{thm:nl} have exponent $2\ell$ and $\ell$. The gap between $\sqrt\ell$ and $\ell$ is open (Open Problem~\ref{op:eth}). Since $w$ is a polynomial in $k$ on the produced instances, no statement in terms of $\ell+w$ jointly is claimed beyond Corollary~\ref{cor:hardeth}.

\emph{Matching the XP algorithm.} Theorem~\ref{thm:hard} shows that the exponent of $n$ in Theorems~\ref{thm:A} and~\ref{thm:nl} cannot be made independent of $\ell$, even though $w$ and $c$ are bounded by functions of $\ell$ alone on the hard instances. Up to the exact exponent of $n$ ($\Omega(\sqrt\ell)$ under ETH, against $\ell$ resp.\ $2\ell$ proved), the complexity of the pair $(\ell,w)$ is therefore settled.

\emph{What is still open.} (1) The reduction has $c=4p+2$ with $p\ge k^2+5$, so $c$ grows with $\ell$; whether CEaMP is W[1]-hard (or fixed-parameter tractable in $\ell$) when $c$ is bounded by a constant, in particular for $c=2$, is open (Open Problem~\ref{op:boundc}). The reason is local to the proof: $p$ is the size of the cliques into which every chain vertex is blown up, and the argument that no clique is split (Lemma~\ref{lem:hsplit}) needs $p>\ell$. (2) The exact exponent under ETH (Open Problem~\ref{op:eth}). (3) Packings with $w\le1$ with tree components of unbounded size, and soft-triangle-free packings (Open Problems~\ref{op:forest} and~\ref{op:BC}); the instances here have many soft triangles and large $w$ (the width is $\Theta(p^2)$ per chain).

\subsection{Overview of the construction}\label{sec:houtline}
The reduction follows the plan of the tunnel attempt of Appendix~\ref{app:tunnels}, and realises the abstract problem of Theorem~\ref{thm:B} with soft triangles; the obstruction of Theorem~\ref{thm:C} is circumvented because the instances are not soft-triangle-free. We use the language of zero-excess (ZE) solutions: a member $(a,b,c)$ with centre $b$ is in state $L$, $R$ or $M$ if its single edit deletes $ab$, deletes $bc$, or inserts $ac$ (states $1,2,3$ of Appendix~\ref{app:tunnels}).
\paragraph{A selection chain.} A \emph{thin chain} $\mathrm{OH}_m$ ($m\ge7$) is an explicit tight instance on the vertices $s_0,\dots,s_{2m},e_0,\dots,e_m$ with spine members $P_j=(s_{2j},s_{2j+1},s_{2j+2})$, further interior members $G_j,F_j$ that carry the vertices $e_j$, two caps at the ends and four hard edges (Appendix~\ref{app:hardness}). Its ZE partitions are exactly $m$ \emph{positions} $a\in\{0,\dots,m-1\}$, in which the spine members are in states $R^{a}ML^{m-1-a}$ (Lemma~\ref{lem:hchain}). This is a tight instance of bounded $c$ and bounded width in which a ``position'' is selected, so that a cut at any one of arbitrarily many places repairs it (Remark~\ref{rem:hchaincore}).

\paragraph{Thick chains.} Every thin vertex $v$ becomes a \emph{super-vertex} $K_v=\{(v,i):i\in\mathbb Z_p\}$, glued by hard edges, and every thin member $(a,b,c)$ is replaced by the $p^2$ \emph{Latin members} $((a,i),(b,j),(c,i+j))$. These cover every pair between $K_a,K_b$, between $K_b,K_c$ and between $K_a,K_c$ exactly once. All pairs of such a triple of super-vertices are soft; so if $i+j\ne k$, the triple $(a,i),(b,j),(c,k)$ is a soft triangle. The purpose is that a ZE partition cannot split a super-vertex (the glue is hard) and that, for an arbitrary solution of excess at most $\ell$, splitting a super-vertex costs more than $\ell$ hard edits when $p\ge\ell+5$ (Lemma~\ref{lem:hsplit}). Hence the thick chain behaves like the thin chain, while a vertex outside the chain can be attached to a super-vertex at a cost that is independent of the position.

\paragraph{Guests, locks and comparators.} A \emph{guest} is a single new vertex $y$ with a set $W_y$ of super-vertices to which it is soft (through $p/2$ \emph{pad members} $(y,(v,i),(v,k))$ for a perfect matching of $K_v$); all other pairs of $y$ are hard non-edges. A guest can join the cluster of a base vertex $h$ only if the entire cluster of $h$ lies in $W_y$; so the set $W_y$ forbids exactly those positions of the chain at which the cluster of $h$ contains a vertex outside $W_y$ (Lemma~\ref{lem:hprof}). Each chain is divided into $U=1+n(n+1)$ \emph{groups} of two consecutive positions. A \emph{permanent guest} sits in the cluster of a host vertex of each group and, for groups that carry no tile, forbids both positions of the group. For a tile $(a,b)\in S_{ij}$ the chain $A_{ij}$ has a \emph{lock guest} and an \emph{anchor} at the slot group $(a,b)$, and the chain $B_{ij}$ at the slot group $(b,a)$; the two anchors of one tile are joined by a hard edge, a \emph{tunnel}. An uncut tunnel forces the lock guest into the host cluster, hence forbids the group; cutting it (one hard edit, cost one) frees the group. So a chain whose tunnels are all uncut has no position, and a chain with exactly one cut tunnel is in a position of the group of that tunnel; as the same hard edit releases the chains $A_{ij}$ and $B_{ij}$, \emph{it chooses the tile $s_{ij}$ in both orders $(a,b)$ and $(b,a)$ at the price of one edit}. Finally, for consecutive chains of a row (of columns) and a separator group $\alpha<n$, a \emph{comparator guest} $y$ with a member $(z,y,z')$ ($z,z'$ the permanent guests of the two chains) must join exactly one of the two chains, and can do so at the first chain only if its position is before the separator and at the second only if it is after; this is the clause $(\beta\le\alpha)\vee(\beta'>\alpha)$ on the blocks $\beta,\beta'$ of the two chosen tiles, and the $n-1$ clauses together express $\beta\le\beta'$. This is Proposition~\ref{prop:hTR}.

\paragraph{Soundness by counting.} The completeness direction is Proposition~\ref{prop:hTR}. For soundness, an arbitrary solution $S$ of excess at most $\ell=k^2$ is analysed without assuming its form. By Lemma~\ref{lem:hsplit} no super-vertex is split. Partition the pairs into regions $\Pi_c$, one for each of the $2k^2$ chains; every region needs at least one edit beyond one per member (Lemma~\ref{lem:hregion}), because with all tunnels uncut the chain has no position. A pair of two regions is a \emph{cross-anchor pair} (a pair between the anchor sets of the two chains of one cell); only these pairs can serve two regions, and they are hard. If $q$ of them are edited, the excess is at least $\max(q,2k^2-q)\ge k^2$ (Theorem~\ref{thm:hS}); if it equals $k^2$, then $S$ is a set of exactly one tunnel cut per cell plus one edit per member, and Proposition~\ref{prop:hTR} gives a tiling. Thus the instance has $\OPT\le|H|+k^2$ iff the Grid Tiling instance is solvable. In this argument the restricted model of Appendix~\ref{app:tunnels} is not assumed; it is derived.


\subsection{Conventions and the thin chain}\label{sec:hthin}\label{sec:hconv}
Soft and hard pairs, members, states and zero-excess (ZE) solutions are as in Section~\ref{sec:prelim} and Appendix~\ref{app:tunnels}. A \emph{ZE partition} of an instance $(G,H)$ is a partition of $V(G)$ whose cost is $|H|$; by Lemma~\ref{lem:decomp} and Lemma~\ref{lem:tZE} it is a partition in which no hard pair is edited (hard edges lie inside clusters, hard non-edges join different clusters) and every member has exactly one edited pair. For a member $(a,b,c)$ with centre $b$ the state is $L$ if the single edit deletes $ab$ ($b\sim c$, $a$ apart), $R$ if it deletes $bc$, and $M$ if it inserts $ac$; these are the states $1,2,3$ of Appendix~\ref{app:tunnels}. Write $x\sim y$ for ``$x,y$ are in one cluster''. The \emph{excess} of a solution $S$ is $\mathrm{ex}(S)=|S|-|H|$. Grid Tiling instances are assumed to have all $S_{ij}\ne\emptyset$ and $n\ge2$ (an instance with an empty $S_{ij}$ is mapped to a fixed NO instance; $n=1$ is trivial).

\paragraph{Definition.} Fix $m\ge7$. The \emph{thin chain} $\mathrm{OH}_m$ has the vertices $s_0,\dots,s_{2m}$ and $e_0,\dots,e_m$ and the following members, listed as $(a,b,c)$ with centre $b$:
\begin{itemize}
\item spine members $P_j=(s_{2j},s_{2j+1},s_{2j+2})$, $0\le j\le m-1$;
\item $G_j=(e_j,s_{2j+1},e_{j+2})$ for $2\le j\le m-4$ and $F_j=(e_j,s_{2j+2},e_{j+1})$ for $2\le j\le m-3$;
\item left cap: $(e_0,s_2,e_1)$, $(s_1,e_1,s_3)$, $(e_0,s_4,e_2)$;
\item right cap: $(s_{2m-5},e_{m-3},s_{2m-1})$, $(e_{m-2},s_{2m-2},e_m)$, $(s_{2m-4},e_{m-1},e_m)$, $(s_{2m-3},e_{m-1},s_{2m})$, $(e_{m-2},e_{m-1},s_{2m-1})$, $(s_{2m-3},e_m,s_{2m-1})$.
\end{itemize}
The edges are the pairs $ab,bc$ of the members and the four \emph{hard edges} $s_0e_1$, $e_0s_3$, $s_{2m-3}e_{m-2}$, $s_{2m}e_m$; every other pair of vertices is a hard non-edge. The members have pairwise disjoint pair sets and are induced $P_3$'s (checked by computer, and by the pair-disjointness argument of Lemma~\ref{lem:hpm} for the thick version). The \emph{cell} of a vertex is $\lfloor t/2\rfloor$ for $s_t$ and $t$ for $e_t$; every member has its vertices in at most three consecutive cells, and every soft pair joins vertices in cells at distance at most $2$. A \emph{position} is an integer $a\in\{0,\dots,m-1\}$; a \emph{group} is a set $\{2u,2u+1\}$, $0\le u\le U-1$, where $m=2U$.

\paragraph{The partitions $P_a$.} Put $\varepsilon(0)=e_1$, $\varepsilon(1)=e_0$, $\varepsilon(j)=e_j$ for $2\le j\le m-2$ and $\varepsilon(m-1)=e_m$. For $0\le a\le m-1$ let $P_a$ be the partition of $V(\mathrm{OH}_m)$ whose clusters $C_0,\dots,C_{m-1}$ are
\[
C_j=\begin{cases}\{s_{2j},s_{2j+1}\}&j<a\\ \{s_{2j},s_{2j+1},s_{2j+2}\}&j=a\\ \{s_{2j+1},s_{2j+2}\}&j>a\end{cases}\ \cup\ \{\varepsilon(j)\}\ \ (\text{for }j\le m-2),
\]
together with the following changes at the right end. If $a\le m-2$ then
\[C_{m-1}=\{s_{2m-1},\allowbreak s_{2m},\allowbreak e_{m-1},\allowbreak e_m\};\]
if $a=m-1$ then $C_{m-2}$ contains in addition $e_{m-1}$ and $C_{m-1}=\{s_{2m-2},\allowbreak s_{2m-1},\allowbreak s_{2m},\allowbreak e_m\}$. In $P_a$ the spine members $P_0,\dots,P_{m-1}$ are in the states $R^{a}ML^{m-1-a}$.

\begin{lemma}[Chain lemma]\label{lem:hchain}
For every $m\ge7$ the ZE partitions of $\mathrm{OH}_m$ are exactly $P_0,\dots,P_{m-1}$; they are pairwise distinct.
\end{lemma}
\begin{proof}[Proof (computer-assisted)]
Both parts are finite statements about a family indexed by $m$ and are proved by computer; the argument that makes the finite check cover all $m\ge7$ is as follows, and is written out in Appendix~\ref{app:hardness}. Every constraint of $\mathrm{OH}_m$ (a member, or a hard edge) involves vertices of at most three consecutive cells, so a condition at a cell depends only on the relation of $a$ to the cell, and on the distance to the left or the right end up to a constant; interior windows are isomorphic, left-end windows do not depend on $m$, and right-end windows depend only on the distance to the right end. \textbf{(CA1)} (part (a)): for every $7\le m\le60$ and every $a$, $P_a$ was checked by computer to have cost exactly $|H|$ with one edit per member and no hard pair edited; all window types occur for $m\le20$. \textbf{(CA2)} (part (b)): the ZE partitions are counted by a transfer-matrix dynamic programme over the cells, with six states per interior cell and an explicit integer $6\times6$ matrix $M$, the same at every interior cell, with $\chi_M(x)=x^4(x-1)^2$; hence $N_m$ is an affine function of $m$ for large $m$, and a direct computation gives $N_m=m$ for $7\le m\le40$. Together with the $m$ distinct partitions $P_a$ of part (a) these are all. Thus the chain lemma is established by a finite computation combined with a translation-invariance argument rather than by a closed-form proof: the claim that the transfer matrix is the same at every interior cell of every $m$ rests on the translation invariance of the member list. The chain lemma for $m\ge10$ is therefore computer-assisted; the Lean formalisation covers $m=7,8,9$ only (Section~\ref{sec:lean}).
\end{proof}

\paragraph{Profiles.} Let $m=2U$. The \emph{host} of the group $u$ is $h_u=s_{4u+2}$; it is the last vertex of $P_{2u}$ and the first vertex of $P_{2u+1}$. Write $\varepsilon^-_u=\varepsilon(2u)$ and $\varepsilon^+_u=\varepsilon(2u+1)$; $\varepsilon^-_u\ne\varepsilon^+_u$ for all $u$.
\begin{lemma}[Profile of a host]\label{lem:hprof}
For the group $u$ the positions fall into four classes $A_0=\{a\le2u-1\}$, $A_1=\{2u\}$, $A_2=\{2u+1\}$, $A_3=\{a\ge 2u+2\}$, and for $a\in A_i$ the cluster of $h_u$ in $P_a$ is $\{h_u\}\cup\Pi_i$ with
\[
\begin{gathered}\Pi_0=\{s_{4u+1},\varepsilon^-_u\},\quad \Pi_1=\{s_{4u},s_{4u+1},\varepsilon^-_u\},\\ \Pi_2=\{s_{4u+3},s_{4u+4},\varepsilon^+_u\},\quad \Pi_3=\{s_{4u+3},\varepsilon^+_u\}.\end{gathered}
\]
($A_0=\emptyset$ for $u=0$ and $A_3=\emptyset$ for $u=U-1$.) In particular $\Pi_0\subset\Pi_1$, $\Pi_3\subset\Pi_2$, $s_{4u}\in\Pi_1\setminus\Pi_0$, $s_{4u+4}\in\Pi_2\setminus\Pi_3$, and $\Pi_0,\Pi_1$ are disjoint from $\Pi_2,\Pi_3$.
\end{lemma}
\emph{Proof} in Appendix~\ref{app:hardness}: the cluster of $h_u$ is read off the definition of $P_a$, with a check of the exceptional $\varepsilon$-labels at the two caps by computer \textbf{(CA3)}.

\paragraph{Allowed sets and forbidden positions.} For a set $W$ of thin vertices of one chain with $h_u\in W$ let $F_u(W)=\{a:\ \text{the cluster of }h_u\text{ in }P_a\not\subseteq W\}$ be the set of positions \emph{forbidden} by $W$. Define, for the group $u$,
\[
\begin{gathered}\mathrm{Pad}_u=\{h_u\}\cup\Pi_0\cup\Pi_3,\quad X_u=\{h_u\}\cup\Pi_0,\quad Y_u=\{h_u\}\cup\Pi_3,\\ \mathrm{All}_u=\{h_u\}\cup\Pi_0\cup\Pi_1\cup\Pi_2\cup\Pi_3 .\end{gathered}
\]
\begin{corollary}\label{cor:hforb}
$F_u(\mathrm{Pad}_u)=\{2u,2u+1\}$, $F_u(X_u)=\{a\ge2u\}$, $F_u(Y_u)=\{a\le2u+1\}$ and $F_u(\mathrm{All}_u)=\emptyset$.
\end{corollary}
\begin{proof}
By Lemma~\ref{lem:hprof}: for $\mathrm{Pad}_u$ the clusters at $A_0$ and $A_3$ are inside $\mathrm{Pad}_u$, while at $A_1$ the vertex $s_{4u}\notin\mathrm{Pad}_u$ and at $A_2$ the vertex $s_{4u+4}\notin\mathrm{Pad}_u$. For $X_u$: $A_0$ is allowed, $A_1$ is forbidden by $s_{4u}$, and $A_2,A_3$ by the non-empty sets $\Pi_2,\Pi_3$ disjoint from $\Pi_0$ (if $A_3=\emptyset$ there is nothing to forbid). $Y_u$ is symmetric: $A_3$ is allowed, $A_2$ is forbidden by $s_{4u+4}$, and $A_1,A_0$ by $\Pi_1,\Pi_0$, which are disjoint from $\Pi_3$. $\mathrm{All}_u$ contains every cluster.
\end{proof}

\subsection{Thick chains and guests}\label{sec:hthick}
Let $p\ge6$ be even. A fixed \emph{one-factorisation} $\mathcal F_1,\dots,\mathcal F_{p-1}$ of $K_p$ on the vertex set $\mathbb Z_p$ (round robin) is used throughout.

\paragraph{Thick chain.} The thick version of a thin chain with tag $c$ has the vertices $(v,i)$, $v$ a thin vertex, $i\in\mathbb Z_p$; $K_v=\{(v,i):i\in\mathbb Z_p\}$ is the \emph{super-vertex} of $v$. For each thin member $(a,b,c')$ there are the $p^2$ \emph{Latin members} $((a,i),(b,j),(c',i+j))$ with centre $(b,j)$, $i,j\in\mathbb Z_p$. For each thin hard pair $uv$ all $p^2$ pairs between $K_u$ and $K_v$ are hard, edges or non-edges as $uv$. Inside each $K_v$ all pairs are hard edges (the \emph{glue}) except the pairs used as pad or host pairs below. Every pair between $K_a$ and $K_b$ ($ab$ a pair of a thin member) is covered by exactly one Latin member: the pair $(a,i),(b,j)$ by the member with these two vertices, the pair $(b,j),(c',k)$ by $i=k-j$, and the pair $(a,i),(c',k)$ by $j=k-i$. So the Latin members have pairwise disjoint pair sets, and if $i+j\ne k$ the triple $(a,i),(b,j),(c',k)$ is a \emph{soft triangle}.

\paragraph{Guests.} A \emph{guest} is a single vertex $y$ together with a set $W_y$ of thin vertices (of one or two chains) such that all pairs between $y$ and the super-vertices $K_v$, $v\in W_y$, are soft and all other pairs of $y$ are hard non-edges, apart from the single-vertex members specified below. Realisation: for every $v\in W_y$ a perfect matching $\mathcal M(y,v)$ of $K_p$ is chosen, a different class $\mathcal F_r$ for different guests at the same $v$; for every $\{i,k\}\in\mathcal M(y,v)$ ($i<k$) there is the \emph{pad member} $(y,(v,i),(v,k))$ with centre $(v,i)$ (edges $y(v,i)$ and $(v,i)(v,k)$, non-edge $y(v,k)$). A \emph{permanent guest} $z$ with a \emph{host} $h\in W_z$ has instead, for the matching $\mathcal M(z,h)$, the \emph{host members} $((h,i),z,(h,k))$ with centre $z$ (edges $z(h,i)$, $z(h,k)$, non-edge $(h,i)(h,k)$), and pad members for all $v\in W_z\setminus\{h\}$. The pairs of $K_v$ in a matching used at $v$ are soft; the glue is the complete graph on $K_v$ minus these pairs. All pairs between different super-vertices that are not soft and not thin hard edges are hard non-edges, and so are all pairs between guests (except those in the members $(x,y,z)$ and $(z,y,z')$ below).

\paragraph{Thin-guest model.} A \emph{thin configuration} is a partition $\pi$ of the set of thin vertices, guests, anchors and stubs of the instance such that: (T1) a pair that is a hard edge in the thin picture (a thin hard edge, an uncut tunnel, a pair $(z,h)$ of a permanent guest and its host) lies inside a cluster; every other pair of two elements is split, except the \emph{free pairs} $(y,v)$, $v\in W_y$, and the pairs of the members; (T2) every thin chain member and every single-vertex member $(x,y,z)$, $(z,y,z')$ has exactly one edited pair.

\begin{lemma}[Blow-up]\label{lem:hblow}
Assume that the glue graph of every $K_v$ is connected (Lemma~\ref{lem:hpar}). Then the ZE partitions of a thick instance are exactly the lifts (replace each thin vertex $v$ by $K_v$) of the thin configurations. In particular in a ZE partition (a) every $K_v$ lies in one cluster; (b) a guest $y$ is together with $K_v$ only if $v\in W_y$; (c) a permanent guest is together with its host $K_h$; (d) the thick members of a thin member are all in the state of the thin member.
\end{lemma}
\begin{proof}
Let $\sigma$ be a ZE partition. No hard pair is edited, so the connected glue puts each $K_v$ into one cluster, which is (a). A pair between $K_u$ and $K_v$ that is a thin hard pair is hard of the same type for all $p^2$ pairs, so $K_u\sim K_v$ iff the thin pair is an edge. For a thin member $(a,b,c')$ the $p^2$ Latin members see the same togetherness pattern of the three super-vertices, so they are in the same state, which is a state of the thin member (one edit); this is (d). A pad member $(y,(v,i),(v,k))$ with $(v,i)\sim(v,k)$ has the togetherness patterns $(1,1,1)$ (if $y\sim K_v$) with one edit, inserting $y(v,k)$, and $(0,1,0)$ (if $y\not\sim K_v$) with one edit, deleting $y(v,i)$; so a pad member imposes no condition. If $v\notin W_y$ the pairs $y(v,i)$ are hard non-edges, so $y\not\sim K_v$, which is (b). A host member $((h,i),z,(h,k))$ with $(h,i)\sim(h,k)$ has one edit iff $z\sim(h,i)$ (pattern $(1,1,1)$: insert $(h,i)(h,k)$); if $z$ is apart it has two edits. This is (c). Conversely, the lift of a thin configuration has no edited hard pair (hard pairs between super-vertices and glue edges are respected by (T1), pairs of guests with $K_v$, $v\notin W_y$, are split) and one edit in every thick member, by the same case analysis. Note that in a lift the host pair $(h,i)(h,k)$ is a soft non-edge inside a cluster and is edited, while a pad pair is a soft edge inside a cluster and is not.
\end{proof}

\subsection{The reduction}\label{sec:hred}
Let $(k,n,(S_{ij}))$ be given, $n\ge2$, and let $p$ be the least even integer $\ge k^2+5$.

\paragraph{Chains and groups.} For every cell $(i,j)\in[k]^2$ there are two chains $A_{ij}$ and $B_{ij}$, each a thick copy of $\mathrm{OH}_m$ with $m=2U$, $U=1+n(n+1)$ (so $m\ge14$). For $\alpha\in[n]$ the \emph{block} $\alpha$ has the \emph{slot groups} $u=(\alpha-1)(n+1)+s$, $s\in[n]$, and is followed by the \emph{separator group} $u_\alpha=\alpha(n+1)$; group $0$ is the left end pad, and the last separator group $u_n=U-1$ is the right end pad.

\paragraph{Tiles and tunnels.} A tile $(a,b)\in S_{ij}$ has its \emph{slot} in $A_{ij}$ at the group $(a-1)(n+1)+b$ (block $a$) and in $B_{ij}$ at the group $(b-1)(n+1)+a$ (block $b$). The groups that are slots of tiles are \emph{tile groups}.

\paragraph{Guests.} For every chain $c$ and every group $u$ there is a permanent guest $z_{c,u}$ with host $h_u$; if $u$ is a tile group, $W_z=\mathrm{All}_u$, and in addition there are a \emph{lock guest} $y_{c,u}$ with $W_y=\mathrm{Pad}_u$, a vertex $x_{c,u}$ (the \emph{anchor}) and the single-vertex member $(x_{c,u},y_{c,u},z_{c,u})$ with centre $y_{c,u}$. If $u$ is not a tile group, $W_z=\mathrm{Pad}_u$ (a \emph{permanent lock}). All other pairs of $x,y,z$ are hard non-edges.

\paragraph{Tunnels.} For every cell and every tile $e\in S_{ij}$ the two anchors of $e$ (in $A_{ij}$ and in $B_{ij}$) are joined by a hard edge, the \emph{tunnel} of $e$. The anchors have no other hard edge.

\paragraph{Comparators.} For $\alpha\in[n-1]$, for every row $i$ and $j<k$, let $c_1=A_{ij}$, $c_2=A_{i,j+1}$, and for every column $j$ and $i<k$ let $c_1=B_{ij}$, $c_2=B_{i+1,j}$. In each case there is a \emph{comparator guest} $y$ with $W_y=X_{u_\alpha}(c_1)\cup Y_{u_\alpha}(c_2)$ (sets of thin vertices of the two chains) and the single-vertex member $(z_{c_1,u_\alpha},y,z_{c_2,u_\alpha})$ with centre $y$.

\paragraph{The instance.} $G$ is defined by the members above (edges $ab,bc$ of every member), the hard edges (thin hard edges, glue, tunnels) and no other edges. $H$ is the set of all members: Latin members, host members, pad members, lock members $(x,y,z)$ and comparator members. The excess is $\ell=k^2$. The instance is denoted $I_{\mathrm{tile}}$; for a set $C$ of tunnels, $I_C$ is $I_{\mathrm{tile}}$ with the pairs of $C$ made non-edges (they stay hard). A \emph{stub instance} is the same with some anchors' tunnels ending at stubs (vertices of degree one).

Lemma~\ref{lem:hpar} below shows that $H$ is a modification-disjoint packing of induced $P_3$'s, that the glue graphs are connected and that matchings suffice, so Lemma~\ref{lem:hblow} applies to $I_{\mathrm{tile}}$, to every $I_C$, and to the stub instances used below.

\subsection{Tunnel-release semantics}\label{sec:hTR}
For a chain $c$ and a cell let $\beta(c)\in[n]$ denote the \emph{block} of the tile whose tunnel is cut: $\beta(A_{ij})=a$ and $\beta(B_{ij})=b$ if the tile is $(a,b)$.

\begin{proposition}[Tunnel release]\label{prop:hTR}
Let $t_{ij}=(a_{ij},b_{ij})\in S_{ij}$ be a tile of every cell, $C$ the set of their tunnels, and $\sigma$ a ZE partition of $I_C$. Then in $\sigma$ every chain $c$ is in a position $a_c$ that lies in the group of its cut tunnel, and $a_{ij}\le a_{i,j+1}$, $b_{ij}\le b_{i+1,j}$ for all cells. Conversely, if $(t_{ij})$ satisfies these inequalities, $I_C$ has a ZE partition.
\end{proposition}
\begin{proof}
($\Rightarrow$) By Lemma~\ref{lem:hblow} take the thin configuration. Its restriction to a chain $c$ is a ZE partition of $\mathrm{OH}_m$ (all non-soft pairs inside the chain are hard), so by Lemma~\ref{lem:hchain} it is $P_{a_c}$ for one position $a_c$.

\emph{Permanent guests.} By Lemma~\ref{lem:hblow}(c), $z=z_{c,u}\sim K_{h_u}$, so $z$ is in the cluster of $h_u$ in $P_{a_c}$; for each $v$ in that cluster the pair $zK_v$ is not a hard non-edge, so $v\in W_z$ by (b). Hence $a_c\notin F_u(W_z)$. If $u$ is not a tile group, Corollary~\ref{cor:hforb} gives $a_c\notin\{2u,2u+1\}$.

\emph{Locks.} Let $u$ be a tile group whose tunnel $xx'$ ($x=x_{c,u}$) is not cut. It is a hard edge of $I_C$ with $N_h(x)=\{x'\}$ and $Q(x)\cap Q(x')=\emptyset$ (the soft partners are guests of different chains), so by Lemma~\ref{lem:tlock} the cluster of $x$ is $\{x,x'\}$; in particular $x\not\sim y,z$ for the members $(x,y,z)$. In the member $(x,y,z)$ with centre $y$ this leaves the single edit ``delete $xy$'' only if $y\sim z$ (if $y\not\sim z$ also $yz$ is deleted: two edits). So $y\sim z$ and $y$ is in the cluster of $h_u$; as above $a_c\notin F_u(\mathrm{Pad}_u)=\{2u,2u+1\}$.

Exactly one tunnel of the chain $c$ is cut (the tunnel of $t_{ij}$ serves both $A_{ij}$ and $B_{ij}$). All groups except the group of that tunnel exclude their two positions, and the groups partition $\{0,\dots,m-1\}$; hence $a_c$ lies in the group of the cut tunnel. With the slot indices above, this gives for the group $u_c$ of $c$: $u_c=(a_{ij}-1)(n+1)+b_{ij}$ for $A_{ij}$ and $(b_{ij}-1)(n+1)+a_{ij}$ for $B_{ij}$. Since $s\in[n]$, the group $u_c$ lies in the block $\beta(c)$: $u_c<u_\alpha$ iff $\beta(c)\le\alpha$, and $u_c>u_\alpha$ iff $\beta(c)\ge\alpha+1$; so $a_c\le2u_\alpha-1$ iff $\beta(c)\le\alpha$ and $a_c\ge2u_\alpha+2$ iff $\beta(c)\ge\alpha+1$.

\emph{Comparators.} Let $y$ be the comparator guest of $(c_1,c_2,\alpha)$, $u=u_\alpha$. In the member $(z_1,y,z_2)$ ($z_i=z_{c_i,u}$) the state $M$ would give $z_1\sim z_2$, impossible because $z_1,z_2$ lie in clusters of different chains (a pair of thick vertices of different chains is a hard non-edge). So the state is $L$ ($y\sim z_1$) or $R$ ($y\sim z_2$). In the first case the cluster of $h_u(c_1)$ is inside $W_y$, i.e.\ inside $X_u(c_1)$ (the part of $W_y$ in $c_1$), so $a_{c_1}\notin F_u(X_u)$, i.e.\ $a_{c_1}\le2u-1$, i.e.\ $\beta(c_1)\le\alpha$. In the second case $a_{c_2}\notin F_u(Y_u)$, i.e.\ $a_{c_2}\ge2u+2$, i.e.\ $\beta(c_2)\ge\alpha+1$. So for every $\alpha\in[n-1]$: $\beta(c_1)\le\alpha$ or $\beta(c_2)>\alpha$, which is equivalent to $\beta(c_1)\le\beta(c_2)$ (if $\beta(c_1)>\beta(c_2)$ take $\alpha=\beta(c_2)\le n-1$). For rows, $\beta(A_{ij})=a_{ij}$, for columns $\beta(B_{ij})=b_{ij}$, which are the inequalities of Grid Tiling.

($\Leftarrow$) Put every chain $c$ at the position $a_c=2u_c$, $u_c$ the group of its cut tunnel. Build the thin configuration: the base clusters are those of $P_{a_c}$. \emph{Permanent guests:} $z_{c,u}$ joins the cluster of $h_u$; this is allowed: for a tile group the set is $\mathrm{All}_u$, for the others $a_c\notin\{2u,2u+1\}$ and Corollary~\ref{cor:hforb} applies. \emph{Locks:} for a tile group $u\ne u_c$ (tunnel not cut) the lock guest $y_{c,u}$ joins the cluster of $z_{c,u}$ (allowed, $a_c\notin\{2u,2u+1\}$), and the anchors $x,x'$ at the two ends of that tunnel form the cluster $\{x,x'\}$; the member $(x,y,z)$ has the single edit ``delete $xy$'' (state $L$). For $u=u_c$ the tunnel is cut, $y_{c,u}$ is not joined to anything but its anchor, and the anchor forms the cluster $\{x,y\}$ (state $R$: delete $yz$). The pair $xx'$ is a cut tunnel, a hard non-edge, and the two anchors are in different clusters. \emph{Comparators:} the guest $y$ of $(c_1,c_2,\alpha)$ joins $z_1$ if $\beta(c_1)\le\alpha$ (allowed: then $a_{c_1}\le2u_\alpha-1$, so the cluster of $h_{u_\alpha}$ is inside $X_{u_\alpha}$) and joins $z_2$ otherwise (then $\beta(c_2)>\alpha$ and $a_{c_2}\ge2u_\alpha+2$, inside $Y_{u_\alpha}$). Every guest is in at most one cluster with a base vertex and two guests are never together, except in the pairs $y,z$ of members: guests join base clusters only through a permanent guest; two hosts $h_u,h_{u'}$ are never in one cluster (their positions in the spine differ by at least four); a tile group carries one lock guest and no comparator, and a separator group carries up to two comparator guests on the same $z$ (one as first chain with condition $\beta\le\alpha$, one as second chain with $\beta\ge\alpha+1$), which cannot both join. All thin chain members have one edit since $P_{a_c}$ is ZE, the members $(x,y,z)$ and $(z_1,y,z_2)$ have one edit as described, and every remaining pair is as required in (T1). By Lemma~\ref{lem:hblow} this lifts to a ZE partition of $I_C$.
\end{proof}

\subsection{Soundness and completeness}\label{sec:hsound}
\begin{lemma}[Parameters and validity]\label{lem:hpar}
(a) $H$ is a modification-disjoint packing of induced $P_3$'s of $G$. (b) Let $u(v)$ be the number of matchings used at $K_v$ (pad matchings and the host matching). Then $u(v)\le3$ for all $v$ and $u(v)\le 2$ if $k\le2$; at a vertex lying in four thin members $u(v)\le2$, and at a vertex in at most three thin members $u(v)\le3$. (c) The glue graph of every $K_v$ is connected, and $p-1\ge3$ matchings are available. (d) $c(H)\le 4p+2$ and $|V(G)|=O(k^2n^2p)$, $|H|=O(k^2n^2p^2)$.
\end{lemma}
\emph{Proof} in Appendix~\ref{app:hardness}; it is a hand count from Lemma~\ref{lem:hprof}, whose result $u_{\max}\le3$, $c\le4p+2$ is also confirmed by computer \textbf{(CA4)}.

\begin{lemma}[No super-vertex is split]\label{lem:hsplit}
Let $S$ be a solution of $I_{\mathrm{tile}}$ with $\mathrm{ex}(S)\le k^2$. Then every $K_v$ lies in one cluster of $G\sd S$. Consequently no hard glue pair and no pad pair (a soft edge inside $K_v$; in particular no pad pair of a comparator pad member) is edited. It is \emph{not} true that no pair inside a super-vertex is edited: the host pairs $(h,i)(h,k)$ are soft non-edges inside $K_h$ and are edited (Lemma~\ref{lem:hblow}).
\end{lemma}
\begin{proof}
Suppose $K_v$ meets at least two clusters; let $X\subset K_v$ be its smallest part, $x=|X|\in[1,p/2]$, $Y=K_v\setminus X$. Every pair between $X$ and $Y$ that is an edge of $G$ is deleted. A vertex of $X$ has at most $u(v)\le3$ soft pairs inside $K_v$, so at least $p-x-u(v)$ hard edges to $Y$; altogether at least $\varphi(x)=x(p-x-u(v))$ hard edits. Since $\varphi$ is concave, $\varphi(x)\ge\min\{\varphi(1),\varphi(p/2)\}$. We have $\varphi(1)=p-1-u(v)\ge p-4\ge k^2+1$. For $k\ge3$: $\varphi(p/2)\ge\frac p2(\frac p2-3)\ge\frac{(k^2+5)(k^2-1)}4>k^2$. For $k\le2$ we have $u(v)\le2$ and $p\in\{6,10\}$: $\varphi(p/2)\ge3>1$ resp.\ $\ge15>4$. So there are more than $k^2$ hard edits, but by Lemma~\ref{lem:decomp} the number of hard edits is at most $\mathrm{ex}(S)\le k^2$, a contradiction. Since $K_v$ lies in one cluster, every glue pair and every pad pair (all of them edges inside $K_v$) is an edge inside a cluster, hence not edited.
\end{proof}
(The host pairs $(h,i)(h,k)$ are soft \emph{non}-edges inside $K_h$; they are edited whenever the permanent guest $z$ joins its host, which happens in every ZE partition (state $M$, the single edit of the host member). What the proofs below use is only: each $K_v$ is inside one cluster, so no hard glue pair and no comparator pad pair is edited.)

\paragraph{Regions.} For a chain $c$ let $R_c$ be the set of its thick vertices and of its guests $z_{c,u},y_{c,u}$ and anchors $x_{c,u}$ (not the comparator guests). Let $T_c$ be the set of tunnel partners of anchors of $c$ (they are anchors of the partner chain, $A_{ij}\leftrightarrow B_{ij}$). The sets $R_c$ are pairwise disjoint, and $T_c\subseteq R_{c'}$ for the partner chain $c'$. Let $H_c$ be the set of members inside $R_c$ (Latin, host, pad members of its permanent and lock guests, lock members), $H_{\mathrm{cmp}}$ the set of members that contain a comparator guest (their pad members and the comparator members); $H$ is the disjoint union of the $H_c$ and $H_{\mathrm{cmp}}$. Let $\Pi_c$ be the set of pairs inside $R_c\cup T_c$ except the pairs with both ends in $T_c$. A \emph{cross-anchor pair} is a pair $\{x,x''\}$ of an anchor $x$ of $c$ carrying a tunnel and an anchor $x''$ of the partner chain carrying a tunnel; cross-anchor pairs are hard (tunnels and hard non-edges). A pair lies in two sets $\Pi_c,\Pi_{d}$ iff it is a cross-anchor pair (and then $d$ is the partner of $c$); every other pair lies in at most one $\Pi$. Pairs of $H_{\mathrm{cmp}}$ members other than the pad pairs involve a comparator guest or join $z_{c_1}$ to $z_{c_2}$ and lie in no $\Pi_c$.

\begin{lemma}[Region bound]\label{lem:hregion}
Let $S$ be a solution with $\mathrm{ex}(S)\le k^2$. For a chain $c$ and a member $M$ let $x_M=S\cap\mathrm{pairs}(M)$; let $h_c$ be the number of edits in $\Pi_c$ on pairs that belong to no member of $H_c$, let $E_c=\sum_{M\in H_c}(|x_M|-1)\ge0$ be the \emph{member excess} (members of $H_c$ with two or three edits), and put $r_c=|S\cap\Pi_c|-|H_c|$. Then $r_c=h_c+E_c\ge\max\{1,h_c\}$.
\end{lemma}
\begin{proof}
By Lemma~\ref{lem:hsplit} the pad pairs of comparator pad members (which lie in $\Pi_c$ but belong to members of $H_{\mathrm{cmp}}$, not of $H_c$) are unedited, so the edits in $\Pi_c$ are on pairs of members of $H_c$ or on pairs of no member at all; this gives $|S\cap\Pi_c|=\sum_{M\in H_c}|x_M|+h_c$, i.e.\ $r_c=h_c+E_c$. Each member has at least one edit, so $E_c\ge0$ and $r_c\ge h_c$; if $h_c\ge1$ we are done. Suppose $h_c=0$ and $E_c=0$, i.e.\ $r_c=0$: then every member of $H_c$ has exactly one edit and $S$ has no other edit in $\Pi_c$. In particular no tunnel at $c$ is cut. Two stubs $x_1',x_2'\in T_c$ (partners of anchors $x_1,x_2$ of $c$) are not together: otherwise $x_1\sim x_1'\sim x_2'\sim x_2$ (tunnels unedited), but $x_1x_2$ is an unedited hard non-edge. So the partition restricted to $R_c\cup T_c$ has no edit on pairs with both ends in $T_c$ either, and is a ZE partition of the instance $G_c=G[R_c\cup T_c]$ with the packing $H_c$ (the comparator pad pairs inside $R_c$ are plain hard edges of $G_c$, and the vertices of $T_c$ are stubs: each has as only edge its tunnel). $G_c$ is a stub instance, so Lemma~\ref{lem:hblow} applies. The proof of Proposition~\ref{prop:hTR} ($\Rightarrow$) for the single chain $c$, with all tunnels present, uses only Lemma~\ref{lem:hchain}, Corollary~\ref{cor:hforb}, the permanent guests and the locks, and no comparators; it shows that the position $a_c$ lies outside every group (every group is excluded either by a permanent lock or by an uncut tunnel), a contradiction since the groups cover all positions. Hence $r_c\ge1$ whenever $h_c=0$, and $r_c\ge\max\{1,h_c\}$ in all cases.
\end{proof}

\begin{theorem}[Soundness]\label{thm:hS}
Every solution $S$ of $I_{\mathrm{tile}}$ has $\mathrm{ex}(S)\ge k^2$. If $\mathrm{ex}(S)=k^2$, then $S$ consists of one cut tunnel in every cell and one edit in every member, and nothing else; hence $S\setminus C$ is a ZE partition of $I_C$ with $C$ the set of cut tunnels, and Grid Tiling has a solution.
\end{theorem}
\begin{proof}
Let $\mathrm{ex}(S)\le k^2$ (otherwise there is nothing to show). Let $q$ be the number of edited cross-anchor pairs and $S_{\mathrm{rest}}$ the set of edits outside $\bigcup_c\Pi_c$. Since only cross-anchor pairs lie in two sets, $|S|=\sum_c|S\cap\Pi_c|-q+|S_{\mathrm{rest}}|$. Every member of $H_{\mathrm{cmp}}$ has an edit; its pad pair is unedited because no comparator pad pair is edited (Lemma~\ref{lem:hsplit}), and its other pairs are in no $\Pi_c$; these members are pair-disjoint, hence $|S_{\mathrm{rest}}|\ge|H_{\mathrm{cmp}}|$. With $r_c=|S\cap\Pi_c|-|H_c|$ (Lemma~\ref{lem:hregion}) and $|H|=\sum_c|H_c|+|H_{\mathrm{cmp}}|$,
\[
\begin{gathered}\mathrm{ex}(S)\ =\ \sum_c r_c-q+\bigl(|S_{\mathrm{rest}}|-|H_{\mathrm{cmp}}|\bigr),\\ r_c=h_c+E_c\ge\max\{1,h_c\},\qquad |S_{\mathrm{rest}}|-|H_{\mathrm{cmp}}|\ge0 .\end{gathered}
\]
All three slacks $E_c$, $r_c-\max\{1,h_c\}$ and $|S_{\mathrm{rest}}|-|H_{\mathrm{cmp}}|$ are non-negative, so $\mathrm{ex}(S)\ge\sum_c\max\{1,h_c\}-q$.
Let $Z=\{c:h_c=0\}$ and let $h_{\mathrm{in}}$ be the number of hard edits lying in exactly one $\Pi_c$. A cross-anchor edit counts in two chains, so $\sum_ch_c=h_{\mathrm{in}}+2q$, and $\sum_c\max\{1,h_c\}=|Z|+\sum_{c\notin Z}h_c=|Z|+h_{\mathrm{in}}+2q$. Hence $\mathrm{ex}(S)\ge|Z|+h_{\mathrm{in}}+q$. Moreover, for $c\notin Z$ we have $r_c=h_c+E_c$, so keeping the member excess of these chains gives $\mathrm{ex}(S)\ge|Z|+h_{\mathrm{in}}+q+\sum_{c\notin Z}E_c$ (the sum runs over $c\notin Z$ only: for $c\in Z$ the term $\max\{1,h_c\}=1$ already accounts for $r_c\ge1$ and $E_c$ must not be added again). The number of chains not in $Z$ is at most $\sum_ch_c$, so $|Z|\ge2k^2-h_{\mathrm{in}}-2q$, and therefore $\mathrm{ex}(S)\ge2k^2-q$. On the other hand every edited cross-anchor pair is a hard edit, so $q\le|U|\le\mathrm{ex}(S)$ (Lemma~\ref{lem:decomp}). Thus $\mathrm{ex}(S)\ge\max\{q,2k^2-q\}\ge k^2$.

Let $\mathrm{ex}(S)=k^2$. Then $q=k^2$ (as $\mathrm{ex}\ge q$ and $\mathrm{ex}\ge2k^2-q$ give $q=k^2$), and every inequality above is tight: from $\mathrm{ex}(S)\ge|Z|+h_{\mathrm{in}}+q+\sum_{c\notin Z}E_c$ we get $Z=\emptyset$, $h_{\mathrm{in}}=0$ and $\sum_{c\notin Z}E_c=0$, hence $\sum_cE_c=0$ because $Z=\emptyset$, and from the first display $|S_{\mathrm{rest}}|=|H_{\mathrm{cmp}}|$. Every chain has $h_c\ge1$ and $\sum_ch_c=2q=2k^2$, so $h_c=1$ for all $2k^2$ chains: the only non-member edit at a chain is one cross-anchor pair. All hard edits are these $k^2$ pairs (as $|U|\le\mathrm{ex}=k^2=q$). Since $E_c=0$, all members of every $H_c$ have exactly one edit, and since $|S_{\mathrm{rest}}|=|H_{\mathrm{cmp}}|$ with one edit at least per member, $S_{\mathrm{rest}}$ has exactly one edit per member of $H_{\mathrm{cmp}}$. In a cell $(i,j)$ exactly one cross-anchor pair is edited (each chain has $h_c=1$ and sees only cross pairs of its own cell). It is a tunnel: suppose the edited pair is a non-edge $x_Ax_B$ of anchors of $A_{ij}$ and $B_{ij}$ whose tunnels are $x_Ax_B^{*}$ and $x_A^{*}x_B$ (with $x_B^*\ne x_B$). It is inserted, so $x_A\sim x_B$; the tunnel $x_Ax_B^*$ is an unedited edge, so $x_A\sim x_B^*$; hence two distinct anchors of $B_{ij}$ are together, but their pair is an unedited hard non-edge ($h_{B_{ij}}=1$ is used up by the cross pair), a contradiction. So the edited cross pair is a tunnel of the cell, deleted, and $S$ is the set $C$ of $k^2$ cut tunnels (one per cell) together with one edit per member. By Lemma~\ref{lem:hsplit} the partition is a partition with $K_v$ intact, and it is a ZE partition of $I_C$. Proposition~\ref{prop:hTR} ($\Rightarrow$) gives a tiling: the tiles $t_{ij}$ of the cut tunnels satisfy $a_{ij}\le a_{i,j+1}$ and $b_{ij}\le b_{i+1,j}$.
\end{proof}

\begin{theorem}[Completeness]\label{thm:hC}
If the Grid Tiling instance has a solution $(s_{ij})$, then $\OPT(I_{\mathrm{tile}})=|H|+k^2$.
\end{theorem}
\begin{proof}
Let $C$ be the tunnels of the tiles $s_{ij}$. By Proposition~\ref{prop:hTR} ($\Leftarrow$) the instance $I_C$ has a ZE partition, whose edit set (relative to $G$) is the set of its $|H|$ member edits together with the $k^2$ deleted tunnels. So $\OPT\le|H|+k^2$, and $\OPT\ge|H|+k^2$ by Theorem~\ref{thm:hS}.
\end{proof}

\subsection{Pathwidth, pair-maximality and proof of Theorem~\ref{thm:hard}}\label{sec:hpwpm}
\begin{lemma}[Pathwidth]\label{lem:hpw}
Every component of $I(H)$ is contained in the set of members of the $k$ chains of one row $\{A_{i,1},\dots,A_{i,k}\}$ or one column $\{B_{1,j},\dots,B_{k,j}\}$ together with their guests and comparator guests; so $I(H)$ has at most $2k$ components, and $\mathrm{pw}(I(H))<24kp^2$, independently of $n$.
\end{lemma}
\emph{Proof} in Appendix~\ref{app:hardness} (a level sweep: members at distance two in level cannot share a vertex, and a level has $O(kp^2)$ members).

\begin{lemma}[Pair-maximality]\label{lem:hpm}
$H$ is pair-maximal.
\end{lemma}
\emph{Proof} in Appendix~\ref{app:hardness} (every non-adjacent pair inside a component of the hard-edge graph is soft).

\begin{proof}[Proof of Theorem~\ref{thm:hard}]
The instance $(G,H,k^2)$ is constructed in time polynomial in $k$ and $n$ (Lemma~\ref{lem:hpar}(d)). (i) holds by definition, Lemma~\ref{lem:hpar}(a),(d) and Lemma~\ref{lem:hpm}; (ii) is Lemmas~\ref{lem:hpar}(d) and~\ref{lem:hpw} ($\tw\le\mathrm{pw}$); (iii) is Theorems~\ref{thm:hS} and~\ref{thm:hC} (a solution of size at most $|H|+k^2$ has excess at most $k^2$, hence exactly $k^2$, hence a tiling exists by Theorem~\ref{thm:hS}). The parameters $\ell+\tw(I(H))$, $\ell+\mathrm{pw}(I(H))$ and $\ell+c$ are bounded by $k^2+24k(k^2+6)^2$, which is a function of $k$ only, and the reduction is polynomial: it is a parameterized reduction from Grid Tiling with $\le$ parameterized by $k$, which is W[1]-hard. An algorithm of running time $f(\ell,w)n^{O(1)}$ would decide Grid Tiling in time $f'(k)\,n^{O(1)}$, so $\mathrm{FPT}=\mathrm{W[1]}$. The ETH statement is Corollary~\ref{cor:hardeth}.
\end{proof}

\begin{remark}[The chain object]\label{rem:hchaincore}
Consider one chain with $N$ tile groups, each with a lock guest and an anchor whose tunnel ends at a stub, all other groups being permanent locks. The proof of Proposition~\ref{prop:hTR} ($\Rightarrow$) shows that with all tunnels present there is no ZE partition, and the proof of ($\Leftarrow$) shows that cutting any single tunnel gives one. In the language of Appendix~\ref{app:tunnels} the set of all $N$ tunnels is therefore a core, of unbounded size $N$, while $c\le4p+2$ and the pathwidth do not depend on $N$ (the bound of Lemma~\ref{lem:hpw}). This is the family of tight instances with bounded $c$ and width and cores of unbounded size that the discussion of Appendix~\ref{app:tunnels} asks for; by Corollary~\ref{cor:sfcore} it must have soft triangles, and it does (the Latin members). The argument is the proof above applied to a single chain. The hypotheses $p$ even, $p\ge8$ suffice for this statement (the tunnel semantics do not use the no-split lemma); for the full problem $p>\ell$ is needed.
\end{remark}

\subsection{Computer assistance}\label{sec:hardcomp}
The proof of Theorem~\ref{thm:hard} is \emph{computer-assisted}. Precisely, the following finite statements are established by computation and are used as lemmas; all the other steps (Lemmas~\ref{lem:hblow}, \ref{lem:hsplit}, \ref{lem:hregion}, Proposition~\ref{prop:hTR}, Theorems~\ref{thm:hS} and~\ref{thm:hC}, the pathwidth bound, the pair-maximality) are proved by hand from them.
\begin{enumerate}
\item[(CA1)] \emph{Lemma~\ref{lem:hchain}(a).} The explicit partitions $P_a$ of the thin chain are ZE partitions for all $m\ge7$ and $a$. Each constraint involves at most three consecutive cells; windows at the left end do not depend on $m$, windows at the right end depend only on the distance to the right end, and interior windows are isomorphic, so a check for $7\le m\le60$ and every $a$ (no failure) covers all $m$.
\item[(CA2)] \emph{Lemma~\ref{lem:hchain}(b).} The chain has no other ZE partition. The number $N_m$ of ZE partitions is computed by a transfer-matrix dynamic programme over the cells (six states per interior cell; an explicit integer matrix of size $6\times6$, identical for all interior cells, with characteristic polynomial $x^4(x-1)^2$); this makes $N_m$ an affine function of $m$ for large $m$, and a computation gives $N_m=m$ for $7\le m\le40$. The statement ``the transfer matrix is the same at every interior cell of every $m$'' rests on translation invariance of the member list and was checked at all cells $7,\dots,19$ of $m=26$.
\item[(CA3)] \emph{Lemma~\ref{lem:hprof}.} The partner sets of the host vertices (the clusters of the hosts at the four classes of positions). The statement is read off the formula of the partitions $P_a$; it was in addition checked by computer for all even $m\in[8,60]$ and all groups ($1{,}836$ cases), because the two cap groups have exceptional $e$-labels.
\item[(CA4)] \emph{Counts.} The values $c\le4p+2$ and $u_{\max}\le3$ (Lemma~\ref{lem:hpar}) are obtained by a count that is easily checked by hand from the sets of Lemma~\ref{lem:hprof}; the count was also carried out by computer for $k\le4$.
\end{enumerate}
Further computational experiments (independent implementations of the chain lemma, tests of the tunnel-release semantics on random Grid Tiling instances, exact optima of full thick instances for $k\le3$) support the construction but are not used in the proof; they are summarised in Appendix~\ref{app:checks}.

\subsection{Where soft triangles are used}
Without soft triangles every instance satisfies the conclusions of Theorem~\ref{thm:sfA} (Section~\ref{sec:soft}): lock cores have at most two anchors, and tunnel releases cannot realise the abstract hardness mechanism (Theorem~\ref{thm:C}). A single chain with its anchors has a core of unbounded size (Remark~\ref{rem:hchaincore}), so by Lemma~\ref{lem:tmaj} its set of ZE state vectors is closed under no majority operation. The Latin members are the soft triangles: the reduction shows that soft triangles suffice for hardness.

\subsection{Consequences for the rest of the paper}
Theorem~\ref{thm:hard} does not contradict any result below: Theorem~\ref{thm:K} needs components of bounded size, but $I(H)$ here has components of size growing with $n$; Theorem~\ref{thm:omega} is fixed-parameter tractable in $(\ell,D,\omega)$ where $D=2c=O(k^2)$ is bounded by a function of $k$, so $\omega$ is unbounded on these instances (the tunnels are hard pairs between far vertices and make the graph $\Gamma'$ of Appendix~\ref{sec:nl} have unbounded treewidth; compare Proposition~\ref{prop:expander}); Theorem~\ref{thm:sfA} is about soft-triangle-free packings. Two consequences follow from results proved earlier. By Proposition~\ref{prop:rpLOC}, the problem of locating a repacking (Open Problem~\ref{op:LOC}) has no fixed-parameter algorithm unless $\mathrm{FPT}=\mathrm{W[1]}$; and a reduction lemma in the sense of Open Problem~\ref{op:ML} cannot exist unless $\mathrm{FPT}=\mathrm{W[1]}$, since by Theorem~\ref{thm:omega} it would give fixed-parameter tractability in $\ell+w$.
